\subsubsection{Interfaces con el usuario}
{\color{black}
El comandante o coordinador interactuará con una aplicación de escritorio para consultar y visualizar escenarios de propagación, cuyos resultados apoyarán la comprensión de la emergencia y la toma de decisiones. Para ello se contemplan los siguientes medios y condiciones:
\begin{itemize}
\item \textbf{Ratón.} Será suficiente un ratón funcional que permita desplazar el puntero y señalar el punto de ignición o el borde activo del incendio, preferiblemente con sensibilidad ajustable para facilitar el trazado de detalles.
\item \textbf{Teclado.} Será suficiente un teclado que permita ingresar valores numéricos en el formulario cuando el servicio meteorológico no esté disponible.
\item \textbf{Pantalla e interfaz gráfica.} La pantalla permitirá visualizar la información y los resultados, mientras que la interfaz se diseñará con PySide6 6.12 para una resolución de 1920 × 1080 píxeles.
\item \textbf{Unidad gráfica.} Se prevé una unidad gráfica integrada en el procesador o dedicada, con 64 MB de memoria RAM, dado que no se esperan altas exigencias de visualización.
\item \textbf{Conexión de red.} Se prevé que una consulta de condiciones meteorológicas a Open-Meteo transfiera alrededor de 5 kB, por lo que una conexión Ethernet o inalámbrica de 64 kbps podría ser suficiente para transmitir esa cantidad en aproximadamente 0,6 segundos bajo condiciones ideales.
\end{itemize}
}
